\section{Shifted moments and positive mixtures}
\label{sec:mixtures}

Equation~\eqref{eq:gaussian-energy} compares the configuration with the
triangular lattice for every kernel $e^{-\pi\alpha t}$.  To pass to a
completely monotone interaction, we represent it by positive mixtures of
these kernels.  A potential may be singular at zero, so a finite representing
measure need not be available for the unshifted function.  We instead
construct one for each positive shift and remove the shift only after taking
the lattice sum.

\subsection{A finite moment representation}

The representation belongs to the classical Hausdorff--Bernstein--Widder
theory \citep{Hausdorff1921,Bernstein1929,Widder1931}.  We prove precisely the
finite-measure form needed here.  The construction first solves a discrete
moment problem using finite differences, then passes to all positive real
exponents.  It also identifies the endpoint masses that enter the energy
argument.

\begin{lemma}[Positive mixtures]\label{lem:positive-mixture}
Let $F:[0,\infty)\to[0,\infty)$ be smooth, with one-sided derivatives at
zero, and suppose $(-1)^jF^{(j)}(t)\ge0$ for every integer $j\ge0$ and every
$t\ge0$.  There is a finite positive Borel measure $\rho$ on $[0,1]$ of
total mass $F(0)$ such that
\begin{equation}\label{eq:positive-mixture}
 F(t)=\int_{[0,1]}v^t\,\dd\rho(v)\qquad(t>0).
\end{equation}
Moreover, $\rho(\{0\})=0$.  An atom at $1$ is allowed and contributes a
constant function.
\end{lemma}

\begin{proof}
\emph{Compactness of the measures.}
We first record an elementary compactness argument that preserves the limiting
total mass without discarding possible endpoint atoms.  Let $(\mu_n)$ be any
sequence of positive measures on $[0,1]$ with
uniformly bounded masses.  By a diagonal subsequence, arrange that the masses
converge to $m$ and that $\mu_n([0,q])$ converges to $a(q)$ for every rational
$q\in(0,1)$.  For $0\le x<1$, set
\[
 H(x)=\inf_{\substack{x<q<1\\q\text{ rational}}}a(q),
 \qquad H(1)=m.
\]
The function $H$ is nondecreasing and right-continuous.  Its associated
positive Borel measure $\mu$ has distribution $\mu([0,x])=H(x)$, including
mass $H(0)$ at zero and mass $m-H(1-)$ at one.  At every continuity point
$x\in(0,1)$ of $H$, rational points on either side of $x$ bracket
$\mu_n([0,x])$ and give convergence to $H(x)$.  Step-function approximations
on partitions with such internal endpoints, together with convergence of
the total masses and uniform continuity, then give
\[
 \int\varphi\,\dd\mu_n\longrightarrow\int\varphi\,\dd\mu
 \quad\text{for every continuous }\varphi:[0,1]\to\R.
\]
Thus the subsequence converges weakly.  The endpoint atoms have not been
discarded in this construction.

\emph{Discrete moments.}
Fix $d>0$ and put $m_j=F(jd)$ for integers $j\ge0$.
Writing $\Delta m_j=m_{j+1}-m_j$, repeated integration gives, for $r\ge1$,
\[
 (-1)^r\Delta^r m_j
 =\int_{[0,d]^r}(-1)^r F^{(r)}(jd+u_1+\cdots+u_r)
       \,\dd u_1\cdots\dd u_r\ge0.
\]
Define a linear functional on polynomials by $\Lambda(x^j)=m_j$.
For every integer $M\ge1$, place at $p/M$, $0\le p\le M$, the mass
\[
 w_{M,p}=\Lambda\!\left(\binom Mp x^p(1-x)^{M-p}\right)
 =\binom Mp(-1)^{M-p}\Delta^{M-p}m_p\ge0.
\]
Let $\mu_{M,d}=\sum_{p=0}^M w_{M,p}\delta_{p/M}$.  Its total mass is
$m_0$, because the Bernstein basis sums to one.  For fixed
$0\le j\le M$, the identity
$\binom Mp\binom pj=\binom Mj\binom{M-j}{p-j}$ gives
\[
 \sum_{p=0}^M\frac{\binom pj}{\binom Mj}
       \binom Mp x^p(1-x)^{M-p}=x^j.
\]
Applying $\Lambda$ shows that
$\sum_p w_{M,p}\binom pj/\binom Mj=m_j$.
For fixed $j$ and $M\ge\max\{1,j\}$, put
\[
 P_{M,j}(x)=\prod_{r=0}^{j-1}\frac{Mx-r}{M-r}.
\]
These polynomials converge uniformly to $x^j$ on $[0,1]$ as $M\to\infty$:
each of their finitely many factors converges uniformly to $x$.  At
$x=p/M$ their values are exactly $\binom pj/\binom Mj$, so the identity
above says $\int P_{M,j}\,\dd\mu_{M,d}=m_j$.  Since $\mu_{M,d}$ is
positive with mass $m_0$,
\[
 \left|\int x^j\,\dd\mu_{M,d}-m_j\right|
 \le m_0\|x^j-P_{M,j}\|_\infty\longrightarrow0.
\]
The compactness argument supplies one weakly convergent subsequence of
$(\mu_{M,d})_{M\ge1}$.  Since the moment convergence holds for every
fixed $j$ along the full sequence, its limit $\eta_d$ satisfies all the moment identities
\[
 F(jd)=\int_{[0,1]}x^j\,\dd\eta_d(x)
 \qquad(j=0,1,2,\ldots).
\]
In particular, $\eta_d$ has mass $F(0)$.

\emph{Positive real exponents.}
For each positive integer $l$, push $\eta_{1/l}$ forward by $x\mapsto x^l$
and call the resulting measure $\rho_l$.  It has mass $F(0)$ and
\[
 \int_{[0,1]}v^{j/l}\,\dd\rho_l(v)=F(j/l)
 \qquad(j=1,2,\ldots).
\]
Choose one weakly convergent subsequence $\rho_{l_n}\to\rho$; necessarily
$l_n\to\infty$, and the limit still has mass $F(0)$.
Now fix any $t>0$ and set $t_n=\lceil l_nt\rceil/l_n$.  Since
$0\le t_n-t<1/l_n$ and
$\sup_{0<v\le1}v^t|\log v|=1/(et)$, the mean value theorem in the exponent
gives the following bound; at $v=0$ both powers are zero because their
exponents are positive:
\[
 \sup_{0\le v\le1}|v^{t_n}-v^t|
 \le\frac{t_n-t}{et}\longrightarrow0.
\]
Continuity of $F$ and weak convergence against the continuous function $v^t$
now imply
\[
 F(t)=\lim_nF(t_n)
 =\lim_n\int v^{t_n}\,\dd\rho_{l_n}
 =\int v^t\,\dd\rho.
\]
The subsequence was chosen before $t$, so the same $\rho$ works for every
$t>0$.  Finally, as $t\downarrow0$, the functions $v^t$ increase to
$1_{(0,1]}(v)$.  Monotone convergence and continuity of $F$ at zero give
$\rho((0,1])=F(0)=\rho([0,1])$, proving $\rho(\{0\})=0$.
\end{proof}

\subsection{Lower limits and singular interactions}

The next proposition isolates the positive-mixture step.  Its Gaussian
hypothesis has already been proved for every configuration considered here
in \eqref{eq:gaussian-energy}.

\begin{proposition}[Positive-mixture transfer]\label{prop:energy-reduction}
Let $\mathcal C\subset\R^2$ be locally finite with centered disk density one.
Suppose that, for every $\alpha>0$,
\[
 \liminf_{R\to\infty}\frac1{N_R}
 \sum_{\substack{x,y\in\mathcal C_R\\x\ne y}}e^{-\pi\alpha|x-y|^2}
 \ge\sum_{a\in A\setminus\{0\}}e^{-\pi\alpha|a|^2}.
\]
Then every smooth completely monotone $g:(0,\infty)\to[0,\infty)$ satisfies
\begin{equation}\label{eq:full-energy-bound}
 E_g(\mathcal C)\ge\sum_{a\in A\setminus\{0\}}g(|a|^2)
\end{equation}
in the extended nonnegative reals, with the ordered-pair energy defined in
\eqref{eq:energy-definition}.
\end{proposition}

\begin{proof}
Fix $\varepsilon>0$ and set $F_\varepsilon(t)=g(\varepsilon+t)$ for $t\ge0$.
It satisfies Lemma~\ref{lem:positive-mixture}, with a finite representing
measure $\rho_\varepsilon$ of mass $g(\varepsilon)$.
Choose any real sequence $R_n\to\infty$, discarding finitely many terms so
that $N_{R_n}>0$, and define
\[
 A_n(v)=\frac1{N_{R_n}}
 \sum_{\substack{x,y\in\mathcal C_{R_n}\\x\ne y}}v^{|x-y|^2}
 \qquad(0\le v\le1).
\]
Every exponent is positive, so $A_n$ is continuous and nonnegative on
$[0,1]$, with $0^t=0$ for $t>0$.  Define the nonnegative extended-valued
function
\[
 L(v)=\sum_{a\in A\setminus\{0\}}v^{|a|^2}\qquad(0\le v\le1).
\]
For $0<v<1$, the Gaussian hypothesis with $\alpha=-\log(v)/\pi$ gives
\[
 \liminf_n A_n(v)\ge L(v).
\]
A full-radius lower-limit inequality holds along every sequence tending to
infinity.  The same comparison holds at the endpoints:
$A_n(0)=L(0)=0$, while
$A_n(1)=N_{R_n}-1\to\infty=L(1)$.

The finite-sum representation, Fatou's lemma, and Tonelli's theorem now give
\begin{align*}
 &\liminf_n\frac1{N_{R_n}}
   \sum_{\substack{x,y\in\mathcal C_{R_n}\\x\ne y}}
        F_\varepsilon(|x-y|^2)\\
 &\qquad=\liminf_n\int A_n(v)\,\dd\rho_\varepsilon(v)\\
 &\qquad\ge\int\liminf_n A_n(v)\,\dd\rho_\varepsilon(v)
 \ge\int L(v)\,\dd\rho_\varepsilon(v)\\
 &\qquad=\sum_{a\in A\setminus\{0\}}F_\varepsilon(|a|^2).
\end{align*}
Every integrand and summand is nonnegative, so these steps remain valid when
the last sum is infinite.  Since the sequence $R_n$ was arbitrary, this is
the full-radius lower bound for $F_\varepsilon$.

Complete monotonicity gives $g'\le0$, hence $g(t)\ge F_\varepsilon(t)$ for
$t>0$.  Termwise comparison therefore yields
\[
 E_g(\mathcal C)\ge
 \sum_{a\in A\setminus\{0\}}g(\varepsilon+|a|^2)
 \qquad(\varepsilon>0).
\]
As $\varepsilon\downarrow0$, the summands increase to $g(|a|^2)$.
Monotone convergence in the lattice sum proves \eqref{eq:full-energy-bound}.
No limit in $\varepsilon$ has been interchanged with the configuration lower
limit, and no unshifted representing measure has been used.
\end{proof}

\begin{proof}[Completion of the proof of Theorem~\ref{thm:universal}]
Equation~\eqref{eq:gaussian-energy} verifies the proposition's hypothesis
for every locally finite configuration of centered disk density one.  This
proves the inequality in Theorem~\ref{thm:universal}.  We finish by checking
attainment, including the case of an infinite lattice sum.

Choose a bounded half-open fundamental cell of $A$ with vertex zero and area one, and let
$d$ bound the distance of its points from its lattice vertex.  Up to boundary
sets of area zero, for $R>d$ the union of cells whose vertices lie in $A\cap B_R$
contains $B_{R-d}$ and is contained in $B_{R+d}$.  Its area is
$\#(A\cap B_R)$, so $A$ has centered disk density one.  For each
$x\in A\cap B_R$, the differences $y-x$ with
$y\in A\cap B_R\setminus\{x\}$ form a subset of $A\setminus\{0\}$.
Nonnegativity gives
\[
 \sum_{\substack{y\in A\cap B_R\\y\ne x}}g(|x-y|^2)
 \le\sum_{a\in A\setminus\{0\}}g(|a|^2).
\]
If the lattice sum is finite, averaging gives the upper bound for $E_g(A)$,
and \eqref{eq:full-energy-bound} applied to $A$ gives the reverse bound.
If the lattice sum is infinite, the same lower bound gives $E_g(A)=\infty$.
Thus the equality in Theorem~\ref{thm:universal} holds in every case.
\end{proof}
